%% kapa: exam=1 round=37 year=2026 period=1 form=A subject=civil_law subject_ko=민법 seq=14 number=14
%% source: Q-Net 감정평가사 기출문제 (cst00309) artlSeq=5258924 file=raw/1008/5258924/2026년 제37회 감정평가사 1차 1교시.pdf
%% conversion: kordoc 4.24.51 -> tools/gen1.py (auto); review=proofread-sonnet; flags=none
\begin{question}{14}
민법상 임의대리에 관한 설명으로 옳은 것은? (다툼이 있으면 판례에 따름)
\begin{choices}
  \item 피한정후견인은 임의대리인이 될 수 없다.
  \item 대리권의 범위가 불분명한 경우, 대리인은 본인을 대리하여 본인의 채권에 대한 소멸시효 중단행위를 할 수 없다.
  \item 법무사가 등기권리자와 등기의무자 쌍방을 대리하여 등기를 신청하는 것은 쌍방대리 금지행위에 해당하지 않는다.
  \item 복임권 없는 대리인이 선임한 복대리인의 권한은 민법 제126조의 권한을 넘은 표현 대리의 기본대리권이 될 수 없다.
  \item 대리인의 능력에 따라 사업의 성공 여부가 결정되는 사무에 관하여 대리권을 수여받은 자라고 하더라도 본인의 묵시적 승낙이 있으면 복대리인을 선임할 수 있다.
\end{choices}
\end{question}
\answer{3}
